\chapter{Lectura proyectada}

\section{Lectura primitiva y lectura derivada}
\begin{definition}
La lectura primitiva de un estado direccional conserva la representación por canales
\[
p=(p_{+x},p_{-x},p_{+y},p_{-y},p_{+z},p_{-z}).
\]
\end{definition}

\begin{definition}
La lectura proyectada es una representación derivada obtenida por diferencia entre canales opuestos.
\end{definition}

\begin{observation}
La lectura proyectada no reemplaza la estructura primitiva. Su función es ofrecer una reducción compatible cuando se desea comparar direcciones mediante coordenadas derivadas.
\end{observation}

\section{Proyección canónica}
\begin{definition}
Dado un estado direccional \(p\), se define su lectura proyectada por
\[
x:=p_{+x}-p_{-x},\qquad
y:=p_{+y}-p_{-y},\qquad
z:=p_{+z}-p_{-z}.
\]
\end{definition}

\begin{observation}
Estas magnitudes deben leerse como coordenadas derivadas. No constituyen la forma originaria del estado dentro de la base carlesina.
\end{observation}

\section{Compatibilidad estructural}
\begin{definition}
La lectura primitiva y la lectura proyectada son compatibles si la segunda puede obtenerse a partir de la primera sin ambigüedad de interpretación dentro del dominio declarado.
\end{definition}

\begin{observation}
La compatibilidad no implica equivalencia total entre ambas lecturas. La estructura por canales conserva información que la proyección puede comprimir.
\end{observation}

\section{Memoria de orientación}
\begin{observation}
La lectura primitiva conserva memoria de orientación porque distingue tránsito por canales opuestos aun cuando la lectura proyectada produzca una misma coordenada derivada final.
\end{observation}

\begin{observation}
Por esta razón, la proyección se admite como herramienta derivada de lectura y no como sustitución completa de la estructura primitiva.
\end{observation}

\section{Compatibilidad con trayectoria}
\begin{definition}
Si una trayectoria mínima se expresa por estados direccionales, su lectura proyectada se obtiene aplicando la proyección a cada uno de sus estados.
\end{definition}

\begin{observation}
La trayectoria proyectada es derivada respecto de la trayectoria primitiva. Su utilidad depende de la conservación suficiente de la información requerida para el análisis local.
\end{observation}

\section{Cláusula de alcance}
\begin{observation}
La lectura proyectada fija una forma derivada de comparación y reducción. No autoriza por sí sola a identificar la proyección con la totalidad de la estructura originaria.
\end{observation}

